\chapter{Estado del arte}

\section{El melanoma: biología y clasificación molecular}

El melanoma es la neoplasia maligna de mayor mortalidad entre los tumores cutáneos. Su incidencia ha aumentado de forma sostenida, especialmente en poblaciones caucásicas con exposición solar acumulada \cite{Potrony2015}. Aunque en estadios localizados puede tratarse quirúrgicamente, el melanoma metastásico presenta un comportamiento clínico agresivo y, hasta la aparición de las terapias modernas, la mediana de supervivencia global era de 6--9 meses \cite{Lebbe2014}.

Desde el punto de vista molecular, el melanoma no es una enfermedad unitaria sino un conjunto de subtipos biológicamente distintos \cite{Bastian2014}. La clasificación más aceptada distingue cuatro subtipos principales según el patrón de mutaciones somáticas: mutante en \textit{BRAF}, mutante en \textit{NRAS}, mutante en \textit{NF1} y triple salvaje. El subtipo \textit{BRAF}$^{V600E}$ es el más frecuente (aproximadamente el 50\,\% de los melanomas cutáneos) y se caracteriza por una activación constitutiva de la vía MAPK, lo que hace que sea el principal objetivo de las terapias dirigidas con inhibidores de BRAF/MEK \cite{Bastian2014, Switzer2022}. Los subtipos \textit{NRAS} y \textit{NF1} también activan la vía MAPK por mecanismos alternativos, mientras que los melanomas triple salvaje son molecularmente más heterogéneos \cite{Podlipnik2021}.

Además de las mutaciones somáticas, existe una predisposición genética relevante. Los genes de alta penetrancia para melanoma incluyen \textit{CDKN2A} (presente en alrededor del 20\,\% de las familias con melanoma), \textit{CDK4}, \textit{BAP1} y \textit{TERT}, mientras que los genes de penetrancia moderada incluyen \textit{MC1R} y \textit{MITF} \cite{Potrony2015}. La carga mutacional tumoral (\textit{tumor mutational burden}, TMB) es un factor adicional de relevancia, ya que se ha asociado al beneficio clínico de los inhibidores de checkpoint inmunológico (ICI) en varios estudios \cite{Podlipnik2021}.

\section{Inhibidores de checkpoint inmunológico: mecanismos de acción}

Los inhibidores de checkpoint inmunológico constituyen la principal estrategia de inmunoterapia en melanoma avanzado. Su fundamento biológico reside en el papel regulador de las moléculas CTLA-4 y PD-1 en el mantenimiento de la tolerancia periférica de los linfocitos T \cite{FifeBluestone2008}.

CTLA-4 (CD152) es un receptor inhibidor expresado en la superficie de las células T que compite con CD28 por la unión a los ligandos B7-1 y B7-2 en las células presentadoras de antígeno. La unión de CTLA-4 reduce la amplitud de la respuesta T en las fases iniciales de activación en el ganglio linfático. PD-1 (CD279), por su parte, actúa principalmente en el lugar tumoral: al unirse a sus ligandos PD-L1 (B7-H1) y PD-L2 (B7-DC), suprime la función efectora de las células T que reconocen el tumor \cite{FifeBluestone2008}. Esta diferencia en el sitio de acción ---nódulo linfático para CTLA-4 y microambiente tumoral para PD-1--- tiene consecuencias sobre el perfil de toxicidad y la biología de la respuesta \cite{Dempke2017}.

Los anticuerpos monoclonales que bloquean estos receptores eliminan los frenos inmunológicos, permitiendo que el sistema inmune reconozca y elimine las células tumorales. Ipilimumab (anti-CTLA-4) fue el primero en demostrar beneficio de supervivencia en melanoma avanzado. Los anti-PD-1, especialmente pembrolizumab y nivolumab, han mostrado mayor eficacia con menor toxicidad y se han convertido en el estándar de primera línea \cite{Switzer2022}.

\section{Evidencia clínica de los inhibidores de checkpoint en melanoma}

\subsection{Anti-PD-1: pembrolizumab y nivolumab}

Hamid et al.\ demostraron por primera vez la actividad de pembrolizumab (entonces denominado lambrolizumab o MK-3475) en melanoma avanzado en un ensayo de fase I con 135 pacientes \cite{Hamid2013}. La tasa de respuesta confirmada fue del 38\,\% en el conjunto de dosis, llegando al 52\,\% en el grupo que recibió 10\,mg/kg cada 2 semanas. Las respuestas fueron duraderas (mediana de 11 meses de seguimiento) y el perfil de toxicidad fue predominantemente de grado 1--2.

Robert et al.\ extendieron estos hallazgos en el estudio KEYNOTE-001, evaluando pembrolizumab en pacientes con melanoma avanzado refractario a ipilimumab \cite{Robert2014}. En esta cohorte de 173 pacientes, la tasa de respuesta global fue del 26\,\% tanto a 2\,mg/kg como a 10\,mg/kg cada 3 semanas, con un perfil de toxicidad manejable. Estos datos establecieron pembrolizumab como una opción eficaz incluso en pacientes previamente tratados con anti-CTLA-4.

\subsection{Anti-CTLA-4: ipilimumab}

Ipilimumab fue aprobado en 2011 tras demostrar beneficio de supervivencia en dos ensayos de fase III. Un análisis agrupado de datos de supervivencia global de 10 estudios prospectivos y dos retrospectivos, incluyendo 1.861 pacientes, mostró una mediana de supervivencia global de 11,4 meses (IC 95\,\%: 10,7--12,1) y una tasa de supervivencia a 3 años del 22\,\%, con una meseta que se inicia alrededor del tercer año de seguimiento y se mantiene hasta los 10 años de observación \cite{Schadendorf2015}. Esta meseta de supervivencia representa uno de los hallazgos más relevantes en el campo y constituye la base para hablar de respuestas duraderas.

Con un seguimiento de 5--6 años, Lebb\'e et al.\ confirmaron que la supervivencia a 5 años con ipilimumab alcanza entre el 12,3\,\% y el 28,4\,\% según la dosis y el estado de tratamiento previo \cite{Lebbe2014}. Este seguimiento prolongado refuerza el concepto de beneficio a largo plazo en un subgrupo de pacientes.

\subsection{Nuevas combinaciones y resistencia}

A pesar de los avances, solo el 25\,\% de los pacientes responde de forma sostenida a la inmunoterapia, y los mecanismos de resistencia primaria y secundaria representan un reto clínico y biológico. Entre los mecanismos descritos se encuentran el agotamiento de células T, la sobreexpresión de caspasa-8 y $\beta$-catenina, la amplificación de los genes PD-L1/PD-L2 y las mutaciones en MHC-I/II \cite{Dempke2017}. Para superar estas resistencias, se evalúan moléculas de segunda y tercera generación dirigidas a nuevas dianas, como TIM-3, LAG-3, VISTA, IDO y KIR, así como anticuerpos coestimuladores como anti-CD40, anti-GITR, anti-OX40 y anti-ICOS \cite{Dempke2017}.

Kim et al.\ exploraron la combinación de ceralasertib (un inhibidor de ATR, enzima clave en la respuesta al daño en el ADN) con durvalumab (anti-PD-L1) en 30 pacientes con melanoma refractario a anti-PD-1 \cite{Kim2022}. La tasa de respuesta objetiva fue del 31\,\% y la tasa de control de la enfermedad del 63,3\,\%, con una duración mediana de respuesta de 8,8 meses. Este estudio subraya el potencial de combinar inmunoterapia con agentes que modulan la reparación del ADN para superar la resistencia.

Una visión de conjunto del panorama terapéutico actual, incluyendo quimioterapia, terapias dirigidas BRAF/MEK, inmunoterapia y terapias combinadas, se puede encontrar en las revisiones de Switzer et al.\ \cite{Switzer2022} y Dhanyamraju y Patel \cite{Dhanyamraju2022}.

\section{Biomarcadores transcriptómicos: pronóstico y predicción}

La identificación de biomarcadores moleculares capaces de estratificar pacientes antes del tratamiento es uno de los objetivos centrales de la oncología de precisión. Se distinguen dos tipos: los biomarcadores pronósticos, que informan sobre la evolución de la enfermedad independientemente de la terapia, y los biomarcadores predictivos, que se asocian a la probabilidad de respuesta a un tratamiento específico.

TCGA-SKCM constituye la cohorte de referencia para el estudio pronóstico en melanoma, al reunir datos de expresión génica, variantes somáticas, metilación y datos clínicos de supervivencia en centenares de pacientes \cite{DasTCGA2020, TCGA}. Los marcadores genéticos de peor pronóstico identificados incluyen mutaciones en \textit{CDKN2A}, amplificaciones de \textit{TERT} y alta carga de daño UV, entre otros \cite{Podlipnik2021}.

En cuanto a la predicción de respuesta a ICI, los factores más estudiados son la expresión de PD-L1, la TMB, la presencia de células T infiltrantes y la firma de IFN-$\gamma$ \cite{Hamid2013, Robert2014}. La transcriptómica de alto rendimiento permite estudiar estos marcadores de forma global y generar firmas multigénicas con potencial predictivo. Sin embargo, su extrapolación entre plataformas (RNA-seq vs.\ NanoString) y tejidos (tumor vs.\ sangre periférica) constituye un reto técnico importante que este trabajo aborda de forma explícita.

\section{Integración de cohortes ómicas}

La disponibilidad de repositorios públicos como GEO y TCGA ha facilitado el análisis de múltiples cohortes con datos de expresión génica y respuesta clínica \cite{GEO, DasTCGA2020}. Sin embargo, la integración directa de cohortes procedentes de plataformas distintas (RNA-seq y NanoString), tejidos diferentes (tumor y sangre periférica) y tratamientos heterogéneos (anti-PD-1 y anti-CTLA-4) introduce riesgos metodológicos relevantes: mezcla de efectos biológicos, restricción al subconjunto de genes comunes y pérdida de potencia estadística.

Por ello, el diseño de este trabajo adopta una arquitectura de descubrimiento, validación y sensibilidad que separa explícitamente las cohortes según su tejido, plataforma y tipo de tratamiento. Esta estrategia, coherente con las recomendaciones metodológicas de la literatura, permite interpretar cada resultado en su contexto biológico y técnico apropiado.

\section{Hipótesis de trabajo}

A partir del estado del arte descrito, la hipótesis central de este proyecto es que los pacientes con melanoma avanzado respondedores a inmunoterapia presentan una señal transcriptómica diferencial relacionada con la activación del sistema inmune, la presentación antigénica por MHC-I y la señalización de IFN-$\gamma$. Esta señal, identificable en el transcriptoma tumoral completo mediante RNA-seq, puede ser evaluada ---con limitaciones derivadas de la distancia biológica y técnica--- en cohortes de validación con otras plataformas, tejidos o tipos de checkpoint.
